% ---------------------------------------------------------------------
%  3  EXPERIMENTAL SETUP
% ---------------------------------------------------------------------
\section{Experimental Setup}
\label{sec:setup}

\subsection{Tools and Environment}
\label{sec:tools}
All code is Python 3.12 with NumPy, pandas, SciPy,
scikit-learn~\cite{pedregosa2011}, PyTorch~\cite{paszke2019}, the
\texttt{timestral} package~\cite{timestral} and Matplotlib; versions are
pinned in \texttt{requirements.txt}. The work consists of five Jupyter
notebooks (base data, data analysis, classical models, neural networks,
Timestral) sharing a \texttt{utils} package for the split dates,
features, metrics and statistical tests. The classical notebook runs in
about 45 minutes on 16 CPU cores, the neural notebook in 33 minutes and
the Timestral notebook in under two hours on a laptop GPU (RTX 5060).
Code: \url{https://github.com/Descalabreitor/IAI600-final-project}.

\subsection{Data Split}
\label{sec:split}
The data is split into three chronological blocks by the date of the
predicted hours: \emph{train} Jan 2020 -- Dec 2022 (26{,}304 hours),
\emph{validation} Jan 2023 -- Jun 2024 (13{,}128) and \emph{test}
Jul 2024 -- Oct 2026 (19{,}899). A 192-hour purge before each border
(the 168-hour LSTM window plus the 24 target hours) ensures that no
target hour of one block lies in the next, so no model is fitted on hours
it is later scored on~\cite{lopezdeprado2018}; inputs that reach further
back are legitimate, since they are known at forecast time. Train and
validation use every sixth origin; the test block uses every origin, and
daily statistics use one forecast per day at 23:00 (829 independent test
days).

\subsection{Model Selection and Hyperparameter Tuning}
\label{sec:tuning}
Every choice (model, hyperparameters, features, refit schedule,
calibration, Timestral setting) is made on train and validation only; the
chosen model is refitted on both and scored once on test. Monthly refits
and rolling calibration are walk-forward: during the test block they use
an outcome only after it is known, as a forecaster in real time would.
Hourly targets are tuned on the validation block, daily targets with four
walk-forward folds over train and validation, which show the true spread
between periods ($\pm0.08$ $R^2$). Grids are small and extended when the
best value lies on an edge. The number of boosting trees is tuned
directly: early stopping on a random hold-out used up to 2.4 times too
many. Models are selected by the reported metric, which matters for
returns: a model trained for the mean but chosen by MAE is pushed towards
the zero forecast. Networks stop early on the last six months of train,
never on a random hold-out, and are averaged over five seeds.

\subsection{Evaluation Metrics and Baselines}
\label{sec:metrics}
Returns are scored by the MAE skill against the zero forecast
($1-\mathrm{MAE}/\mathrm{MAE}_0$), by $\Roos$ (the $R^2$ against the zero
forecast) and by the information coefficient (IC, the correlation of
forecast and outcome). Direction is scored by the AUC and by the accuracy
of the 10\% most confident calls (binomial test); size and volatility by
$\Rtwo$; ranges by coverage, width and CRPS~\cite{gneiting2007}. The
$\Rtwo$ of $|\ret|$ has a low ceiling: a single hour's $|\ret|$ is a very
noisy measure of the volatility behind it, so even a perfect volatility
forecast explains only part of its variance, and values of 0.1--0.2
already mean real skill. Baselines are the zero forecast, the reversal of
the last hour, last week's volatility times the hour profile, and
``always up''. Uncertainty is measured by a moving-block
bootstrap~\cite{kunsch1989} over the test days (blocks of 10 days, 1{,}000
resamples) and Diebold--Mariano tests~\cite{diebold1995}; a difference is
called \emph{reliable} when its 95\% interval excludes zero.
